\documentclass[12pt,a4paper]{article}
\usepackage[utf8]{inputenc}
\usepackage[T2A]{fontenc}
\usepackage[russian]{babel}
\usepackage{graphicx}
\usepackage{booktabs}
\usepackage{geometry}
\usepackage{float}
\usepackage{paracol}
\usepackage{src/title}
\usepackage{hyperref}
\usepackage{amsmath}
\usepackage{listings}
\usepackage{xcolor}

\geometry{top=2cm, left=2cm, right=2cm, bottom=2cm}
\emergencystretch=1.5em
\sloppy

\title{Лабораторная работа №6. Логистическая регрессия --- теория и вопросы защиты}
\author{Ануфриев А.С. 465029}
\supervisor{Болдырева Елена Александровна}
\date{08.11.2026}

\abstract{
Теоретические сведения по логистической регрессии и ответы на вероятные вопросы защиты лабораторной работы: математическая форма модели, функция потерь, алгоритмы обучения, метрики классификации и исследование гиперпараметров на примере датасета пассажиров Титаника.
}
\keywords{логистическая регрессия, сигмоида, log-loss, градиентный спуск, метод Ньютона, метрики классификации}

\begin{document}
\maketitle

\section{Теоретические сведения}

\subsection{Постановка задачи}
Логистическая регрессия предназначена для бинарной классификации. Каждому объекту ставится в соответствие метка \(y \in \{0,1\}\). Модель оценивает вероятность принадлежности объекта к положительному классу:

\begin{equation}
P(y=1 \mid x) = \sigma(z), \qquad z = x^T w,
\end{equation}

где \(w\) --- вектор параметров модели (включая свободный член), \(x\) --- вектор признаков объекта, \(\sigma(\cdot)\) --- сигмоидная функция.

\subsection{Сигмоидная функция}
\begin{equation}
\sigma(z) = \frac{1}{1 + e^{-z}}, \qquad \sigma'(z) = \sigma(z)(1-\sigma(z)).
\end{equation}
Функция \(\sigma\) монотонно возрастает, её значения лежат на отрезке \((0,1)\), что позволяет интерпретировать выход модели как вероятность. Порог классификации обычно принимается равным 0.5: объект относится к классу 1, если \(p \ge 0.5\), и к классу 0 в противном случае.

\subsection{Функция потерь}
Для обучения используется функция логарифмических потерь (log-loss, кросс-энтропия):

\begin{equation}
L(w) = -\frac{1}{N}\sum_{i=1}^{N}\Bigl[y_i\ln p_i + (1-y_i)\ln(1-p_i)\Bigr],
\end{equation}

где \(p_i = \sigma(x_i^T w)\) --- предсказанная вероятность для \(i\)-го объекта. Функция выпукла по параметрам \(w\), поэтому градиентные методы сходятся к глобальному минимуму.

\subsection{Градиентный спуск}
Градиент функции потерь по параметрам:

\begin{equation}
\nabla_w L = \frac{1}{N}X^T(p-y),
\end{equation}

где \(X\) --- матрица признаков со столбцом смещения, \(p\) и \(y\) --- векторы предсказанных вероятностей и истинных меток. Обновление весов:

\begin{equation}
w \leftarrow w - \eta \nabla_w L,
\end{equation}

где \(\eta\) --- скорость обучения (learning rate). Выбор \(\eta\) влияет на скорость сходимости и стабильность: слишком большое \(\eta\) может привести к дивергенции, слишком маленькое --- к медленной сходимости.

\subsection{Оптимизация Ньютона}
Вместо градиентного спуска можно использовать метод Ньютона (второй порядок). Гессиан функции потерь:

\begin{equation}
H = \frac{1}{N}X^T W X,
\end{equation}

где \(W\) --- диагональная матрица с элементами \(p_i(1-p_i)\) на диагонали. Обновление:

\begin{equation}
w \leftarrow w - H^{-1}\nabla_w L.
\end{equation}

Метод Ньютона обычно сходится за меньшее число итераций, но требует обращения (или решения линейной системы) с матрицей размером \((d+1) \times (d+1)\), где \(d\) --- число признаков, что дороже по вычислениям на каждой итерации.

\subsection{Метрики классификации}
Построим матрицу ошибок на тестовой выборке:
\begin{itemize}
\item TP --- истинно положительные (выжили, предсказано выжили);
\item TN --- истинно отрицательные (не выжили, предсказано не выжили);
\item FP --- ложноположительные (не выжили, предсказано выжили);
\item FN --- ложноотрицательные (выжили, предсказано не выжили).
\end{itemize}

Тогда:
\begin{equation}
\text{Accuracy} = \frac{TP+TN}{TP+TN+FP+FN},
\end{equation}
\begin{equation}
\text{Precision} = \frac{TP}{TP+FP}, \qquad
\text{Recall} = \frac{TP}{TP+FN}, \qquad
\text{F1-Score} = \frac{2 \cdot \text{Precision} \cdot \text{Recall}}{\text{Precision} + \text{Recall}}.
\end{equation}

Accuracy измеряет общую долю верных ответов. Precision --- долю верно предсказанных выживших среди всех, кому модель предсказала выживание. Recall --- долю реально выживших, которые модель верно обнаружила. F1-Score --- гармоническое среднее precision и recall, устойчивое к désbalansсу классов.

\section{Вероятные вопросы защиты с ответами}

\subsection*{Вопрос 1. В чём отличие логистической регрессии от линейной?}
\textbf{Ответ.} Линейная регрессия предсказывает непрерывное числовое значение и minimизирует среднеквадратичную ошибку. Логистическая регрессия предназначена для классификации: через сигмоиду она выдаёт вероятность принадлежности к классу в интервале \((0,1)\), а функцией потерь служит кросс-энтропия (log-loss). Порог 0.5 превращает вероятность в класс 0 или 1.

\subsection*{Вопрос 2. Зачем нужна сигмоида и какова её производная?}
\textbf{Ответ.} Сигмоида отображает любое действительное число \(z\) на отрезок \((0,1)\), что даёт интерпретацию выхода как вероятности. Её производная \(\sigma'(z) = \sigma(z)(1-\sigma(z))\) используется при вычислении градиента функции потерь и при построении гессиана для метода Ньютона.

\subsection*{Вопрос 3. Как выбирается порог классификации, и что изменится при его сдвиге?}
\textbf{Ответ.} Стандартный порог --- 0.5. Повышение порога снижает число ложноположительных ответов (растёт precision) и увеличивает число ложноотрицательных (падает recall). Понижение порога действует наоборот. Выбор порога зависит от задачи: например, при скрининге важнее recall, при целевой рассылке --- precision.

\subsection*{Вопрос 4. Как обрабатывались пропуски в датасете Титаника?}
\textbf{Ответ.} Пропуски в возрасте (\emph{Age}) и стоимости билета (\emph{Fare}) заполнены медианными значениями, пропуски в порту посадки (\emph{Embarked}) --- модальным. Признак \emph{Cabin} с большим числом пропусков исключён из модели. Категориальные признаки закодированы числовыми значениями, числовые признаки стандартизированы.

\subsection*{Вопрос 5. Какое влияние оказывает скорость обучения на обучение модели?}
\textbf{Ответ.} Маленькая скорость обучения обеспечивает стабильную, но медленную сходимость. Большая скорость ускоряет сходимость на начальных этапах, но может вызывать колебания весов и дивергенцию, особенно на поздних итерациях. В работе наилучший результат получен при learning rate 0.01; при lr = 0.1 и 0.5 качество на тесте не улучшалось и даже снижалось при длительном обучении.

\subsection*{Вопрос 6. Зачем нужна стандартизация признаков в логистической регрессии?}
\textbf{Ответ.} Признаки имеют разные размерности и диапазоны (например, возраст до 80 лет, стоимость билета до 512). Стандартизация приводит их к сравнимому масштабу, что ускоряет сходимость градиентного спуска и делает шаги обновления весов сбалансированными по всем признакам.

\subsection*{Вопрос 7. В чём преимущество и недостаток метода Ньютона по сравнению с градиентным спуском?}
\textbf{Ответ.} Метод Ньютона использует вторую производную и обычно требует меньше итераций для сходимости. Недостаток --- необходимость обращать матрицу гессиана размера \((d+1)\times(d+1)\) на каждой итерации, что существенно дороже при большом числе признаков. В работе метод Ньютона сходился быстро, но не превзошёл градиентный спуск по точности на тестовой выборке.

\subsection*{Вопрос 8. Как оценивалось качество модели и какая метрика выбрана главной?}
\textbf{Ответ.} Качество оценивалось на отдельной тестовой выборке (20\% данных) по accuracy, precision, recall и F1-Score. Главной метрикой принят F1-Score, так как он учитывает и precision, и recall и устойчив к désbalansсу классов (выживших меньше, чем не выживших). Наилучший F1-Score 0.7402 получен при градиентном спуске с lr = 0.01 и 1000 итерациями.

\subsection*{Вопрос 9. Что означают значения TP, TN, FP, FN для лучшей модели?}
\textbf{Ответ.} Для лучшей комбинации на тестовой выборке: TP = 47 (верно предсказано выживших), TN = 99 (верно предсказано не выживших), FP = 15 (ошибочно предсказано выжившими), FN = 18 (выжившие, ошибочно отнесённые к не выжившим). На основании них вычислены accuracy 0.8156, precision 0.7581, recall 0.7231, F1-Score 0.7402.

\subsection*{Вопрос 10. Можно ли сказать, что логистическая регрессия линейная модель?}
\textbf{Ответ.} Да, логистическая регрессия --- линейная модель в том смысле, что решение о классе принимается по линейной комбинации признаков \(z = x^T w\): сигмоида монотонна, поэтому \(p \ge 0.5\) эквивалентно \(x^T w \ge 0\), что задаёт линейную границу разделения классов в пространстве признаков.

\end{document}
